% NOTE (2026-08-30): theory appendix created in the audit-driven restructure.
% Contains the proofs for Section 4 (sec:theory). Style rules: never use the
% em-dash; upright (non-italic) body for theorem-like environments.
\section{Proofs for Section~\ref{sec:theory}}
\label{app:theory}

Throughout this appendix the standing assumptions of the setup paragraph of
Section~\ref{sec:theory} are in force. Because every metric $H\succ0$ makes
$\langle u,Hv\rangle$ an inner product on a finite-dimensional space, all
statements are statements about a Hilbert space, and we work directly in
$\|\cdot\|_H$ without rescaling.

\subsection{Convergence and rate for certified iterations}

\begin{proposition}[Krasnosel'ski\u{\i}--Mann convergence with a
last-iterate rate]
\label{prop:km}
Let $H\succ0$, let $T$ be $\theta$-averaged in $\|\cdot\|_H$ with
$\theta\in(0,1)$, and let $\operatorname{Fix}T\neq\emptyset$. Then the
iteration $z_{k+1}=Tz_k$ converges, from any $z_0$, to some
$z^\star\in\operatorname{Fix}T$; the residual $\|z_{k+1}-z_k\|_H$ is
nonincreasing in $k$; and, for every $k\geq0$,
\begin{equation}
 \|z_{k+1}-z_k\|_H^2
 \;\leq\;
 \frac{\theta}{(1-\theta)\,(k+1)}\;
 \operatorname{dist}_H^2\!\left(z_0,\operatorname{Fix}T\right).
 \label{eq:km-rate}
\end{equation}
\end{proposition}

\begin{proof}
Write $T=(1-\theta)I+\theta S$ with $S$ nonexpansive in $\|\cdot\|_H$; all
norms in this proof are $\|\cdot\|_H$. Fix any
$z^\star\in\operatorname{Fix}T$. Expanding $z_{k+1}=(1-\theta)z_k+\theta
Sz_k$ with the identity
$\|(1-\theta)a+\theta b\|^2=(1-\theta)\|a\|^2+\theta\|b\|^2
-\theta(1-\theta)\|a-b\|^2$, at $a=z_k-z^\star$ and $b=Sz_k-z^\star$, and
bounding $\|Sz_k-z^\star\|\leq\|z_k-z^\star\|$ by nonexpansiveness gives the
Fej\'er inequality
\begin{equation}
 \|z_{k+1}-z^\star\|^2
 \;\leq\;
 \|z_k-z^\star\|^2-\frac{1-\theta}{\theta}\,\|z_{k+1}-z_k\|^2 .
 \label{eq:fejer}
\end{equation}
The map $T$ is nonexpansive, so
$\|z_{k+2}-z_{k+1}\|\leq\|z_{k+1}-z_k\|$. Summing \eqref{eq:fejer} over the
first $k+1$ iterations, keeping only the last (smallest) residual, and
minimizing over $z^\star$ yields \eqref{eq:km-rate}. By \eqref{eq:fejer} the
iterates are bounded and $\|z_{k+1}-z_k\|\to0$, so by continuity every
cluster point of $(z_k)$ is fixed by $T$; reusing \eqref{eq:fejer} with such
a cluster point as $z^\star$ shows $\|z_k-z^\star\|$ is nonincreasing and
vanishes along a subsequence, hence $z_k\to z^\star$. This argument is the
Krasnosel'ski\u{\i}--Mann theorem with its modern rate
\citep[Theorem~1]{RyuYin2022LargeScale}; see also
\citet{Krasnoselskii1955Two}, \citet{Mann1953Mean}, and
\citet{BauschkeCombettes2017Convex}.
\end{proof}

\begin{remark}[Monitored residual, normalization, and summability]
\label{rem:monitored-residual}
The theorem controls $\|z_{k+1}-z_k\|_{H_G}$, the step residual of the
certified map that is actually iterated. For a genome
$\mathrm{relax}_\rho(G')$ the compiled map is
$T_G=(1-\rho)I+\rho T_{G'}$, so
$z_{k+1}-z_k=\rho\,(T_{G'}z_k-z_k)$: the two residuals differ by the factor
$\rho$, which matters when comparing genomes with different relaxation genes,
and \eqref{eq:km-rate} divided by $\rho^2$ bounds the unrelaxed fixed-point
residual $\|T_{G'}z_k-z_k\|_{H_G}^2$. Summing \eqref{eq:fejer} also gives the
summability estimate
$\sum_{k=0}^{\infty}\|z_{k+1}-z_k\|_{H_G}^2\leq
\frac{\theta_G}{1-\theta_G}\operatorname{dist}_{H_G}^2
(z_0,\operatorname{Fix}T_G)$.
\end{remark}

\subsection{Base-scheme certificates}

\begin{lemma}[Base schemes]
\label{lem:base-certificates}
For each base scheme of Table~\ref{tab:certificates} with positive stepsizes
satisfying the gate inequality, the pair $(H_G,\theta_G)$ of the table, with
$\theta_G=1/(2-c_G)$, together with the scheme's standard decoder $\pi_G$,
is a construction certificate in the sense of
Definition~\ref{def:certificate}.
\end{lemma}

\begin{proof}
We first establish every row with the true norm $\|L\|$ in place of
$\ell$, then account for the conservative bound.

Since $\nabla f$ is $(1/L_f)$-cocoercive by the Baillon--Haddad theorem
\citep{BaillonHaddad1977Quelques}, the forward step $I-\alpha\nabla f$ is
$(\alpha L_f/2)$-averaged for $\alpha\in(0,2/L_f)$, and the proximal step is
$\tfrac12$-averaged. A composition of $\theta_1$- and $\theta_2$-averaged
maps is $(\theta_1+\theta_2-2\theta_1\theta_2)/(1-\theta_1\theta_2)$-averaged
\citep[Theorem~27]{RyuYin2022LargeScale}, a constant due to
\citet{OguraYamada2002Nonstrictly}; see also
\citet{CombettesYamada2015Compositions}. At $\theta_1=\tfrac12$ and
$\theta_2=c$ it is $1/(2-c)$; this is the FBS row, and the Davis--Yin row is
the same constant \citep[Theorem~28]{RyuYin2022LargeScale}, established by
\citet{DavisYin2017Three}. DRS is
$\tfrac12\bigl(I+R_{\alpha A}R_{\alpha B}\bigr)$, a composition of reflected
resolvents, hence firmly nonexpansive for every $\alpha>0$
\citep[\S2.7]{RyuYin2022LargeScale}, a fact going back to
\citet{LionsMercier1979Splitting}. For the primal--dual rows,
$M_{\tau\sigma}\succ0$ exactly when $\tau\sigma\|L\|^2<1$, and in that metric
PDHG \citep{ChambollePock2011First} is a proximal-point iteration on the
saddle subdifferential, hence firmly nonexpansive in
$\|\cdot\|_{M_{\tau\sigma}}$ \citep[\S3.3]{RyuYin2022LargeScale}; see also
the contraction analysis of \citet{HeYuan2012Convergence}. Condat--V\~u
\citep{Condat2013Primal,Vu2013Splitting} is variable-metric
forward--backward in the same metric \citep{CombettesVu2014Variable}, and its
forward step is $c$-averaged in $\|\cdot\|_{M_{\tau\sigma}}$ with
$c=\tau L_f/\bigl(2(1-\tau\sigma\|L\|^2)\bigr)$, so composition gives
$1/(2-c)$ again \citep[\S3.3]{RyuYin2022LargeScale}. The inequality
$\tau L_f/2+\tau\sigma\|L\|^2<1$ is equivalent, for $\tau,\sigma>0$, to the
conjunction of $\tau\sigma\|L\|^2<1$ and $c<1$, so a single inequality
delivers both (C1) and a positive-definite metric.

Now let $\ell\geq\|L\|$ be the supplied bound and let $c_G$ denote the
tabulated forward-step constant computed from $\ell$. The checked
inequalities in $\ell$ imply the corresponding inequalities in $\|L\|$, so
$H_G\succ0$, and the true forward-step constant $c$ computed from $\|L\|$
satisfies $c\leq c_G<1$, so the compiled iteration is $\theta$-averaged with
$\theta=1/(2-c)$. A $\theta$-averaged map is also $\theta'$-averaged for
every $\theta'\in[\theta,1)$: writing $T=(1-\theta)I+\theta S$ gives
$T=(1-\theta')I+\theta'S'$ with $S'=(1-\theta/\theta')I+(\theta/\theta')S$,
a convex combination of nonexpansive maps. Since
$\theta=1/(2-c)\leq1/(2-c_G)=\theta_G<1$, the pair $(H_G,\theta_G)$
certifies (C1), exactly when $\ell=\|L\|$ and conservatively otherwise.

For (C2) and (C3), no bijection between fixed points and solutions is
needed; it suffices that fixed points exist and that every fixed point
decodes to a solution. For the primal--dual schemes the state is the saddle
pair and, because $M_{\tau\sigma}\succ0$, the fixed points of the
(variable-metric) proximal-point and forward--backward iterations are
exactly the zeros of the saddle operator, i.e.\ the saddle set, which is
nonempty by assumption; $\pi_G$ is the primal coordinate
\citep[\S3.3]{RyuYin2022LargeScale}. For FBS and Davis--Yin the fixed points
are exactly the solutions of the associated inclusion (for Davis--Yin after
applying the decoder $\pi_G=\operatorname{prox}_{\alpha g}$), and for DRS
every fixed point $z$ satisfies
$\operatorname{prox}_{\alpha g}(z)\in\mathcal X^\star$ while every solution
arises from at least one fixed point; these standard correspondences are
\citep[\S2.7]{RyuYin2022LargeScale}. In each case the nonempty saddle set
yields a fixed point through the correspondence, giving (C2), and the
decoder maps all of $\operatorname{Fix}T_G$ into $\mathcal X^\star$, giving
(C3).
\end{proof}

\iclrTODO{Base-scheme instantiation table: for each row of
Table~\ref{tab:certificates}, give the compiled state space, the exact
update the compiler emits for problem~\eqref{eq:composite} (including how
FBS and DRS instantiate the three-term objective), the decoder, and the
fixed-point correspondence, and verify the entries against the
implementation. Without this table, Lemma~\ref{lem:base-certificates}
certifies the textbook operators, not yet the emitted code.}

\subsection{Closure rules}

\begin{lemma}[Relaxation]
\label{lem:relaxation}
Let $T$ be $\theta$-averaged in $\|\cdot\|_H$ and let $\rho>0$ satisfy
$\rho\theta<1$. Then $(1-\rho)I+\rho T$ is $\rho\theta$-averaged in
$\|\cdot\|_H$ with $\operatorname{Fix}\bigl((1-\rho)I+\rho
T\bigr)=\operatorname{Fix}T$.
\end{lemma}

\begin{proof}
Writing $T=(1-\theta)I+\theta S$ with $S$ nonexpansive gives
$(1-\rho)I+\rho T=(1-\rho\theta)I+\rho\theta S$, which is
$\rho\theta$-averaged because $\rho\theta\in(0,1)$, and $\rho>0$ makes the
fixed-point equations equivalent. The decoder and metric are untouched, so
the certificate transports as stated in Theorem~\ref{thm:soundness}(ii).
\end{proof}

% PROPOSED to Nikhil: the operational account of blending (mechanism, the
% semantic fixed-point caveat, and the empirical negative) now lives in
% Appendix app:blending, demoted from the main text; the old iclrTODO asking
% for one consistent account is resolved by that demotion.
The next lemma is the certificate rule for the blend extension discussed in
Appendix~\ref{app:blending}. It is excluded from
Theorem~\ref{thm:soundness} because its hypothesis of a nonempty common
fixed-point set is semantic: the current gate has no syntactic test for it.

\begin{lemma}[Blends]
\label{lem:blend}
Let $T_1,T_2$ be $\theta_1$- and $\theta_2$-averaged in a common
$\|\cdot\|_H$ with
$\operatorname{Fix}T_1\cap\operatorname{Fix}T_2\neq\emptyset$, and let
$\lambda\in(0,1)$. Then $T=\lambda T_1+(1-\lambda)T_2$ is
$\bigl(\lambda\theta_1+(1-\lambda)\theta_2\bigr)$-averaged in $\|\cdot\|_H$
and $\operatorname{Fix}T=\operatorname{Fix}T_1\cap\operatorname{Fix}T_2$.
\end{lemma}

\begin{proof}
Write $T_i=(1-\theta_i)I+\theta_iS_i$ and set
$\theta=\lambda\theta_1+(1-\lambda)\theta_2$. Then
$T=(1-\theta)I+\theta\bigl(\tfrac{\lambda\theta_1}{\theta}S_1
+\tfrac{(1-\lambda)\theta_2}{\theta}S_2\bigr)$, and the bracket is a convex
combination of nonexpansive maps, hence nonexpansive; this is the
combination rule of \citet{CombettesYamada2015Compositions}. The inclusion
$\operatorname{Fix}T\supseteq\operatorname{Fix}T_1\cap\operatorname{Fix}T_2$
is immediate. Conversely, let $z\in\operatorname{Fix}T$ and
$w\in\operatorname{Fix}T_1\cap\operatorname{Fix}T_2$. Then
$\|z-w\|\leq\lambda\|T_1z-w\|+(1-\lambda)\|T_2z-w\|\leq\|z-w\|$, so both
inequalities are equalities and $\|T_iz-w\|=\|z-w\|$ for $i=1,2$; the
Fej\'er inequality \eqref{eq:fejer} applied to $T_i$ at $z$ then forces
$\|T_iz-z\|=0$. Hence $z\in\operatorname{Fix}T_1\cap\operatorname{Fix}T_2$.
\end{proof}

\subsection{The Halpern wrapper}

\begin{proposition}[Halpern genomes]
\label{prop:halpern}
Let $T_G$ carry a certificate $(H_G,\theta_G,\pi_G)$ and set
$S_G=I+\theta_G^{-1}(T_G-I)$. Then $S_G$ is nonexpansive in
$\|\cdot\|_{H_G}$ with $\operatorname{Fix}S_G=\operatorname{Fix}T_G$, and the
compiled iteration
\begin{equation}
 z_{k+1}=\lambda_{k+1}z_{\rm anc}+(1-\lambda_{k+1})S_Gz_k ,
 \label{eq:halpern}
\end{equation}
under any admitted schedule, i.e.\ $\lambda_k\to0$,
$\sum_k\lambda_k=\infty$, and $\sum_k|\lambda_{k+1}-\lambda_k|<\infty$,
converges to $P^{H_G}_{\operatorname{Fix}T_G}(z_{\rm anc})$, the projection
of the anchor onto $\operatorname{Fix}T_G$ in $\|\cdot\|_{H_G}$
\citep{Halpern1967Fixed,Wittmann1992Approximation}. The canonical schedule
$\lambda_{k+1}=1/(k+2)$ meets all three conditions and, when the wrapper
anchors at the initial iterate ($z_{\rm anc}=z_0$), in addition gives
\begin{equation}
 \|S_Gz_k-z_k\|_{H_G}^2
 \;\leq\;
 \frac{4\operatorname{dist}_{H_G}^2\!\left(z_0,\operatorname{Fix}
 T_G\right)}{(k+1)^2}
 \label{eq:halpern-rate}
\end{equation}
\citep{Lieder2021Optimized,Kim2021Accelerated}.
\end{proposition}

\begin{proof}
$S_G$ is precisely the nonexpansive map in the averaged decomposition
$T_G=(1-\theta_G)I+\theta_GS_G$ of (C1), and $T_Gz=z$ if and only if
$S_Gz=z$ because $\theta_G>0$. Since $(\mathcal Z_G,\langle\cdot,H_G
\cdot\rangle)$ is a Hilbert space, the cited Halpern results apply verbatim
in $\|\cdot\|_{H_G}$: convergence to the projection of the anchor is
\citet{Wittmann1992Approximation}, and the rate \eqref{eq:halpern-rate} is
\citet{Lieder2021Optimized}.
\end{proof}

Rate \eqref{eq:halpern-rate} is an $O(1/k^2)$ bound on the same squared
residual that \eqref{eq:km-rate} bounds at $O(1/k)$, and it is exactly
optimal in the following sense: \citet{ParkRyu2022Exact} exhibit a matching
lower bound for deterministic fixed-point methods that access $S_G$ only
through evaluations, with iterates constrained to the span of the anchor and
past evaluations. The claim is optimality within that oracle class, not an
unrestricted impossibility statement.

\subsection{Finite-prefix switching}
% PROPOSED to Nikhil: replaces summable-error safeguard with finite-prefix
% switch, matches implementation

\begin{lemma}[Finite-prefix switching]
\label{lem:switchk}
Let $T_B$ be $\theta$-averaged in $\|\cdot\|_H$ on a state space
$\mathcal Z$ with $\operatorname{Fix}T_B\neq\emptyset$, let
$T_A:\mathcal Z\to\mathcal Z$ be an arbitrary map, and fix
$K\in\mathbb N$. Define $z_{k+1}=T_Az_k$ for $k<K$ and $z_{k+1}=T_Bz_k$
for $k\geq K$. If $z_K$ is finite, then $(z_k)$ converges to some
$z^\star\in\operatorname{Fix}T_B$, and for every $k\geq K$ the residual
bound \eqref{eq:km-rate} holds with $z_0$ replaced by $z_K$ and $k+1$
replaced by $k-K+1$.
\end{lemma}

\begin{proof}
For $k\geq K$ the sequence is exactly the certified iteration
$z_{k+1}=T_Bz_k$ started at the finite point $z_K$.
Proposition~\ref{prop:km} applies to this orbit and gives convergence to a
point of $\operatorname{Fix}T_B$ together with the rate from the new
starting point. The prefix enters only through the constant
$\operatorname{dist}_H(z_K,\operatorname{Fix}T_B)$, so it cannot affect the
asymptotics.
\end{proof}

The compiler enforces the hypotheses: $K$ is a bounded gene, and the
handoff is guarded by a finiteness check on $z_K$. If a prefix produces
a nonfinite state, the implementation starts the tail from the last finite
state, or the initial state if necessary, and records that recovery. This
recovered trajectory is distinct from the unrecovered prefix in the lemma.
The rule deliberately trades any guarantee about the prefix for total
freedom in choosing it: $T_A$ needs no certificate, no error budget, and no
per-step monitor, in contrast to summable-error analyses of inexact
Krasnosel'ski\u{\i}--Mann iterations
\citep{Combettes2001QuasiFejerian,Combettes2004Solving}, inexact
Douglas--Rachford and ADMM \citep{EcksteinBertsekas1992Douglas}, and
safeguarded Anderson acceleration \citep{Zhang2020Globally}, which admit
perturbations at every iteration at the price of a per-step error test.

\FloatBarrier
